\documentclass[11pt, a4paper]{article}\usepackage[]{graphicx}\usepackage[]{xcolor}
% maxwidth is the original width if it is less than linewidth
% otherwise use linewidth (to make sure the graphics do not exceed the margin)
\makeatletter
\def\maxwidth{ %
  \ifdim\Gin@nat@width>\linewidth
    \linewidth
  \else
    \Gin@nat@width
  \fi
}
\makeatother

\definecolor{fgcolor}{rgb}{0.345, 0.345, 0.345}
\newcommand{\hlnum}[1]{\textcolor[rgb]{0.686,0.059,0.569}{#1}}%
\newcommand{\hlsng}[1]{\textcolor[rgb]{0.192,0.494,0.8}{#1}}%
\newcommand{\hlcom}[1]{\textcolor[rgb]{0.678,0.584,0.686}{\textit{#1}}}%
\newcommand{\hlopt}[1]{\textcolor[rgb]{0,0,0}{#1}}%
\newcommand{\hldef}[1]{\textcolor[rgb]{0.345,0.345,0.345}{#1}}%
\newcommand{\hlkwa}[1]{\textcolor[rgb]{0.161,0.373,0.58}{\textbf{#1}}}%
\newcommand{\hlkwb}[1]{\textcolor[rgb]{0.69,0.353,0.396}{#1}}%
\newcommand{\hlkwc}[1]{\textcolor[rgb]{0.333,0.667,0.333}{#1}}%
\newcommand{\hlkwd}[1]{\textcolor[rgb]{0.737,0.353,0.396}{\textbf{#1}}}%
\let\hlipl\hlkwb

\usepackage{framed}
\makeatletter
\newenvironment{kframe}{%
 \def\at@end@of@kframe{}%
 \ifinner\ifhmode%
  \def\at@end@of@kframe{\end{minipage}}%
  \begin{minipage}{\columnwidth}%
 \fi\fi%
 \def\FrameCommand##1{\hskip\@totalleftmargin \hskip-\fboxsep
 \colorbox{shadecolor}{##1}\hskip-\fboxsep
     % There is no \\@totalrightmargin, so:
     \hskip-\linewidth \hskip-\@totalleftmargin \hskip\columnwidth}%
 \MakeFramed {\advance\hsize-\width
   \@totalleftmargin\z@ \linewidth\hsize
   \@setminipage}}%
 {\par\unskip\endMakeFramed%
 \at@end@of@kframe}
\makeatother

\definecolor{shadecolor}{rgb}{.97, .97, .97}
\definecolor{messagecolor}{rgb}{0, 0, 0}
\definecolor{warningcolor}{rgb}{1, 0, 1}
\definecolor{errorcolor}{rgb}{1, 0, 0}
\newenvironment{knitrout}{}{} % an empty environment to be redefined in TeX

\usepackage{alltt}

\usepackage[top = 0.6 in, bottom = 0.6 in, left = 0.8 in, right = 0.8 in]{geometry}
\usepackage{amsmath, amssymb, amsfonts}

\allowdisplaybreaks[4]

\usepackage{enumerate}
\usepackage{array}
\usepackage{multirow}
\usepackage{dingbat}
\usepackage{fontawesome5}
\usepackage{tasks}
\usepackage{bbding}
\usepackage{twemojis}
% how to use bull's eye ----- \scalebox{2.0}{\twemoji{bullseye}}
\usepackage{fontspec}
\usepackage{fancyhdr}
\usepackage{lastpage}

\pagestyle{fancy}
\fancyhf{}
\fancyfoot[C]{Page \thepage\ of \pageref{LastPage}}
\renewcommand{\headrulewidth}{0pt}

\title{MSMS 408 : Practical 02}
\author{{\fontsize{22}{20} \benyorithafont Ananda Biswas} \\[1em] Exam Roll No. : 24419STC053}
\date{\today}

\newfontface\benyorithafont{Benyoritha-G334D.otf}

\newfontface\myfont{Myfont1-Regular.ttf}[LetterSpace=0.05em]

\newfontface\cbfont{CaveatBrush-Regular.ttf}
% how to use --- \myfont --write text here--
\IfFileExists{upquote.sty}{\usepackage{upquote}}{}
\begin{document}

\maketitle


\section*{\faArrowAltCircleRight[regular] \textcolor{blue}{Question}}

\hspace{1cm} Consider a random sample $\{2, 3, 1, 4, 0, 2, 3, 5, 1, 2, 4, 3, 2, 1, 3, 4, 2, 3, 1, 2\}$ drawn from a Poisson distribution with parameter $\lambda$. Assume that the prior distribution of $\lambda$ follows a Gamma distribution with mean $2.5$ and coefficient of variation $0.5$.

\begin{enumerate}[(i)]
\item Determine the parameters of the Gamma distribution.

\item Derive the posterior distribution of $\lambda$.

\item Compute the posterior mean and posterior variance of $\lambda$.

\item Obtain the Bayes estimate of $\lambda$ under the squared error loss function.

\item Generate a sample of size $10000$ from the posterior distribution and plot the prior and posterior densities along with a histogram of the generated posterior sample.

\item Provide a $95\%$ Bayesian credible interval for $\lambda$ based on simulated posterior values. Also obtain the theoretical quantiles and compare the two.
\end{enumerate}

\section*{\faArrowAltCircleRight[regular] \textcolor{blue}{Theory}}

\begin{itemize}

\item[\scalebox{2.0}{\twemoji{keycap: 1}}] \hspace{0.1cm} Let the gamma distribution has shape $ = \alpha$ and rate $= \beta$. Then, mean $ = \dfrac{\alpha}{\beta}$ and variance $ = \dfrac{\alpha}{\beta^2}$. So, c.v. $ = \dfrac{1}{\sqrt{\alpha}}$.\\[0.3em] 

Here $\dfrac{1}{\sqrt{\alpha}} = 0.5 \Rightarrow \alpha = 4$, consequently $\beta = \dfrac{4}{2.5} = 1.6$.\\[0.3em]

So, $\lambda \sim \text{Gamma}(4, 1.6)$.

\item[\scalebox{2.0}{\twemoji{keycap: 2}}] \hspace{0.1cm} The prior distribution is 

$$\lambda \sim g(\lambda) \propto e^{-\beta \lambda} {\lambda}^{\alpha - 1}.$$

The likelihood of $\lambda$ is

$$L(\lambda) \propto e^{-n \lambda} {\lambda}^{\sum \limits_{i = 1}^{n} x_i}.$$

$\therefore$ The posterior distribution of $\lambda$ will be

$$g(\lambda | \mathbf{x}) \propto e^{-\beta \lambda - n \lambda } {\lambda}^{\alpha + \sum \limits_{i = 1}^{n} x_i - 1}.$$

Recall that, Gamma is a conjugate prior for Poisson parameter. So we identify the posterior distribution of $\lambda$ as

$$\lambda | \mathbf{x} \sim \text{Gamma}\left( \alpha + \sum \limits_{i = 1}^{n} x_i, n + \beta \right).$$

\begin{knitrout}
\definecolor{shadecolor}{rgb}{0.969, 0.969, 0.969}\color{fgcolor}\begin{kframe}
\begin{alltt}
\hldef{alpha} \hlkwb{<-} \hlnum{4}\hldef{; beta} \hlkwb{<-} \hlnum{1.6}
\end{alltt}
\end{kframe}
\end{knitrout}

\begin{knitrout}
\definecolor{shadecolor}{rgb}{0.969, 0.969, 0.969}\color{fgcolor}\begin{kframe}
\begin{alltt}
\hldef{x} \hlkwb{<-} \hlkwd{c}\hldef{(}\hlnum{2}\hldef{,} \hlnum{3}\hldef{,} \hlnum{1}\hldef{,} \hlnum{4}\hldef{,} \hlnum{0}\hldef{,} \hlnum{2}\hldef{,} \hlnum{3}\hldef{,} \hlnum{5}\hldef{,} \hlnum{1}\hldef{,} \hlnum{2}\hldef{,} \hlnum{4}\hldef{,} \hlnum{3}\hldef{,} \hlnum{2}\hldef{,} \hlnum{1}\hldef{,} \hlnum{3}\hldef{,} \hlnum{4}\hldef{,} \hlnum{2}\hldef{,} \hlnum{3}\hldef{,} \hlnum{1}\hldef{,} \hlnum{2}\hldef{)}
\end{alltt}
\end{kframe}
\end{knitrout}

From the given, $\lambda | \mathbf{x} \sim \text{Gamma}(52, 21.6)$.

\item[\scalebox{2.0}{\twemoji{keycap: 3}}] \hspace{0.1cm} The posterior mean and variance are

$$E[\lambda | \mathbf{x}] = \dfrac{\alpha + \sum \limits_{i = 1}^{n} x_i}{n + \beta} = 2.41, \qquad \text{Var}[\lambda | \mathbf{x}] = \dfrac{\alpha + \sum \limits_{i = 1}^{n} x_i}{(n + \beta)^2} = 0.111.$$

\item[\scalebox{2.0}{\twemoji{keycap: 4}}] \hspace{0.1cm} Under squared error loss function, Bayes estimate of $\lambda$ is the posterior mean $2.41$.

\end{itemize}

\section*{\faArrowAltCircleRight[regular] \textcolor{blue}{R Program}}





\begin{itemize}
\item[\scalebox{2.0}{\twemoji{keycap: 5}}]
\end{itemize}

\begin{knitrout}
\definecolor{shadecolor}{rgb}{0.969, 0.969, 0.969}\color{fgcolor}\begin{kframe}
\begin{alltt}
\hldef{n} \hlkwb{<-} \hlnum{10}\hlopt{^}\hlnum{5}
\end{alltt}
\end{kframe}
\end{knitrout}

\begin{knitrout}
\definecolor{shadecolor}{rgb}{0.969, 0.969, 0.969}\color{fgcolor}\begin{kframe}
\begin{alltt}
\hldef{lambda_posterior} \hlkwb{<-} \hlkwd{rgamma}\hldef{(n,} \hlkwc{shape} \hldef{=} \hlnum{52}\hldef{,} \hlkwc{rate} \hldef{=} \hlnum{21.6}\hldef{)}
\end{alltt}
\end{kframe}
\end{knitrout}

\begin{knitrout}
\definecolor{shadecolor}{rgb}{0.969, 0.969, 0.969}\color{fgcolor}\begin{kframe}
\begin{alltt}
\hlkwd{data.frame}\hldef{(lambda_posterior)} \hlopt{%>%}
  \hlkwd{ggplot}\hldef{(}\hlkwd{aes}\hldef{(}\hlkwc{x} \hldef{= lambda_posterior))} \hlopt{+}
  \hlkwd{geom_histogram}\hldef{(}\hlkwd{aes}\hldef{(}\hlkwc{y} \hldef{=} \hlkwd{after_stat}\hldef{(density)),} \hlkwc{bins} \hldef{=} \hlnum{20}\hldef{,}
                 \hlkwc{fill} \hldef{=} \hlsng{"#F54927"}\hldef{,} \hlkwc{color} \hldef{=} \hlsng{"black"}\hldef{,} \hlkwc{linewidth} \hldef{=} \hlnum{1}\hldef{)} \hlopt{+}
  \hlkwd{stat_function}\hldef{(}\hlkwc{fun} \hldef{= dgamma,} \hlkwc{args} \hldef{=} \hlkwd{list}\hldef{(}\hlkwc{shape} \hldef{=} \hlnum{4}\hldef{,} \hlkwc{rate} \hldef{=} \hlnum{1.6}\hldef{),} \hlkwc{xlim} \hldef{=} \hlkwd{c}\hldef{(}\hlnum{0}\hldef{,} \hlnum{5.3}\hldef{),}
                \hlkwc{linewidth} \hldef{=} \hlnum{1.25}\hldef{,} \hlkwc{linetype} \hldef{=} \hlsng{"dashed"}\hldef{)} \hlopt{+}
  \hlkwd{stat_function}\hldef{(}\hlkwc{fun} \hldef{= dgamma,} \hlkwc{args} \hldef{=} \hlkwd{list}\hldef{(}\hlkwc{shape} \hldef{=} \hlnum{52}\hldef{,} \hlkwc{rate} \hldef{=} \hlnum{21.6}\hldef{),}
                \hlkwc{linewidth} \hldef{=} \hlnum{1.25}\hldef{)} \hlopt{+}
  \hlkwd{labs}\hldef{(}\hlkwc{x} \hldef{=} \hlkwd{expression}\hldef{(lambda))}
\end{alltt}
\end{kframe}
\includegraphics[width=\maxwidth]{figure/unnamed-chunk-7-1} 
\end{knitrout}

\newpage

\begin{itemize}
\item[\scalebox{2.0}{\twemoji{keycap: 6}}] $95\%$ Bayesian credible interval based on simulated posterior sample is given by $$(q_{0.025}, q_{0.975})$$ where $q_{\alpha}$ is the $\alpha-$fractile of the simulated sample.

\begin{knitrout}
\definecolor{shadecolor}{rgb}{0.969, 0.969, 0.969}\color{fgcolor}\begin{kframe}
\begin{alltt}
\hlkwd{quantile}\hldef{(lambda_posterior,} \hlkwd{c}\hldef{(}\hlnum{0.025}\hldef{,} \hlnum{0.975}\hldef{))}
\end{alltt}
\begin{verbatim}
##     2.5%    97.5% 
## 1.798726 3.110130
\end{verbatim}
\end{kframe}
\end{knitrout}


The credible interval based on theoretical quantiles would be
\begin{knitrout}
\definecolor{shadecolor}{rgb}{0.969, 0.969, 0.969}\color{fgcolor}\begin{kframe}
\begin{alltt}
\hlkwd{qgamma}\hldef{(}\hlkwd{c}\hldef{(}\hlnum{0.025}\hldef{,} \hlnum{0.975}\hldef{),} \hlkwc{shape} \hldef{=} \hlnum{52}\hldef{,} \hlkwc{rate} \hldef{=} \hlnum{21.6}\hldef{)}
\end{alltt}
\begin{verbatim}
## [1] 1.797967 3.104425
\end{verbatim}
\end{kframe}
\end{knitrout}

\end{itemize}

\smallpencil \hspace{0.1cm} {\setlength{\spaceskip}{1em plus 0.5em minus 0.5em} \fontsize{17}{20}\myfont Credible interval obtained from simulation is consistent with the theoretical one. \par}

\end{document}
